\chapter{微积分基础}

\intro{微积分是理解深度学习优化算法的关键。梯度下降法建立在导数和梯度的概念之上，反向传播算法本质上是链式法则的系统应用。本章系统介绍极限与导数、偏导数与梯度、链式法则、Hessian矩阵、Taylor展开、凸优化等核心主题。}

\section{极限与连续性}

\subsection{极限的定义}

\begin{mydefinition}{函数的极限}
设函数 $f(x)$ 在点 $x_0$ 的某个去心邻域内有定义。若存在常数 $L$，使得当 $x$ 无限接近 $x_0$ 时，$f(x)$ 无限接近 $L$，则称 $L$ 为 $x\to x_0$ 时 $f(x)$ 的极限，记为：
\[
\lim_{x\to x_0} f(x) = L
\]
\end{mydefinition}

精确的 $\varepsilon$-$\delta$ 定义：$\forall\varepsilon>0$，$\exists\delta>0$，当 $0<|x-x_0|<\delta$ 时，有 $|f(x)-L|<\varepsilon$。

\subsection{连续性}

若 $\lim_{x\to x_0}f(x)=f(x_0)$，则称 $f$ 在 $x_0$ 处连续。深度学习中的常见激活函数（ReLU、sigmoid、tanh）在定义域上都是连续的，但ReLU在 $x=0$ 处的导数不连续，这在数学上不影响优化（因为梯度下降使用的是次梯度）。

\section{导数}

\subsection{导数的定义}

\begin{mydefinition}{导数}
函数 $f(x)$ 在点 $x$ 处的导数定义为：
\[
f'(x) = \lim_{h\to 0} \frac{f(x+h)-f(x)}{h}
\]
若该极限存在，则称 $f$ 在 $x$ 处可导。
\end{mydefinition}

导数的几何意义是函数曲线在该点处切线的斜率，物理意义是瞬时变化率。在深度学习中，导数衡量了损失函数关于参数的"瞬时"变化率，指示了调整参数的方向——沿着负梯度方向移动可以减小损失函数的值。

\subsection{基本导数公式}

\begin{myformula}
\frac{d}{dx}(c) &= 0 \\
\frac{d}{dx}(x^n) &= n x^{n-1} \\
\frac{d}{dx}(e^x) &= e^x \\
\frac{d}{dx}(\ln x) &= \frac{1}{x} \\
\frac{d}{dx}(\sin x) &= \cos x \\
\frac{d}{dx}(\cos x) &= -\sin x \\
\frac{d}{dx}(a^x) &= a^x \ln a
\end{myformula}

\subsection{导数的运算法则}

\begin{myTheorem}{导数运算法则}
设 $f$ 和 $g$ 在 $x$ 处可导，则：
\begin{itemize}
    \item \textbf{线性性}：$(af+bg)' = af' + bg'$（其中 $a,b$ 为常数）
    \item \textbf{乘法法则}：$(fg)' = f'g + fg'$
    \item \textbf{除法法则}：$\left(\frac{f}{g}\right)' = \frac{f'g-fg'}{g^2}$
    \item \textbf{链式法则}：$(f(g(x)))' = f'(g(x)) \cdot g'(x)$
\end{itemize}
\end{myTheorem}

\subsection{常见激活函数的导数}

深度学习中的激活函数及其导数具有重要地位：

\begin{itemize}
    \item \textbf{Sigmoid}：$\sigma(x)=\frac{1}{1+e^{-x}}$，$\sigma'(x)=\sigma(x)(1-\sigma(x))$
    \item \textbf{Tanh}：$\tanh(x)=\frac{e^x-e^{-x}}{e^x+e^{-x}}$，$\tanh'(x)=1-\tanh^2(x)$
    \item \textbf{ReLU}：$\operatorname{ReLU}(x)=\max(0,x)$，$\operatorname{ReLU}'(x)=\begin{cases}0 & x<0\\1 & x>0\end{cases}$（在 $x=0$ 处导数未定义，实践中通常取 $0$ 或使用次梯度）
    \item \textbf{Leaky ReLU}：$f(x)=\max(\alpha x, x)$，$\alpha$ 为一个小的正数（如 $0.01$）
    \item \textbf{GELU}：高斯误差线性单元，常用于Transformer模型
\end{itemize}

\begin{mycode}{Python中的激活函数与导数}
\begin{lstlisting}[language=Python]
import numpy as np

def sigmoid(x):
    return 1 / (1 + np.exp(-x))

def sigmoid_derivative(x):
    s = sigmoid(x)
    return s * (1 - s)

def relu(x):
    return np.maximum(0, x)

def relu_derivative(x):
    return (x > 0).astype(float)

def tanh_derivative(x):
    t = np.tanh(x)
    return 1 - t ** 2

# 测试
x = np.array([-2.0, -1.0, 0.0, 1.0, 2.0])
print(f"x: {x}")
print(f"sigmoid(x): {sigmoid(x)}")
print(f"sigmoid'(x): {sigmoid_derivative(x)}")
print(f"ReLU(x): {relu(x)}")
print(f"ReLU'(x): {relu_derivative(x)}")
\end{lstlisting}
\end{mycode}

\section{偏导数与梯度}

\subsection{偏导数}

对于多元函数 $f(x_1,x_2,\ldots,x_n)$，关于 $x_i$ 的偏导数定义为在固定其他变量情况下，$f$ 关于 $x_i$ 的变化率：

\[
\frac{\partial f}{\partial x_i} = \lim_{h\to 0} \frac{f(x_1,\ldots,x_i+h,\ldots,x_n)-f(x_1,\ldots,x_i,\ldots,x_n)}{h}
\]

偏导数的计算规则与一元导数相同，只是将其他变量视为常数。

\begin{example}
求 $f(x,y)=x^2y+3xy^2$ 的偏导数。

\[
\frac{\partial f}{\partial x} = 2xy + 3y^2, \quad \frac{\partial f}{\partial y} = x^2 + 6xy
\]
\end{example}

\subsection{梯度}

\begin{mydefinition}{梯度}
函数 $f(x_1,x_2,\ldots,x_n)$ 的梯度是由所有偏导数组成的向量：
\[
\nabla f = \begin{bmatrix}
\frac{\partial f}{\partial x_1} & \frac{\partial f}{\partial x_2} & \cdots & \frac{\partial f}{\partial x_n}
\end{bmatrix}^{\top}
\]
\end{mydefinition}

梯度的几何意义：梯度向量 $\nabla f$ 指向函数值增长最快的方向，其模长表示该方向上的最大变化率。因此，$-\nabla f$ 指向函数值下降最快的方向——这正是梯度下降算法的核心思想。

\begin{figure}[H]
\centering
\begin{tikzpicture}[
    myarrow/.style={->, >=stealth, thick},
    mylabel/.style={font=\small}
]
    % 3D-like contour plot
    \draw[->, >=stealth] (-4, 0) -- (4, 0) node[right] {$x_1$};
    \draw[->, >=stealth] (0, -4) -- (0, 4) node[above] {$x_2$};

    % Contour lines (concentric ellipses)
    \foreach \r in {0.5, 1, 1.5, 2, 2.5, 3, 3.5} {
        \draw[blue!40, thick, rotate around={30:(0,0)}]
            (0, 0) ellipse ({\r*1.5} and \r);
    }

    % Point
    \fill[red] (2.5, 1.5) circle (3pt);
    \node[mylabel, red, right] at (2.5, 1.5) {$\mathbf{x}$};

    % Gradient vector
    \draw[myarrow, red!80] (2.5, 1.5) -- (1.5, 0.5);
    \node[mylabel, red!80] at (1.3, 0.3) {$-\nabla f$};

    % Tangent to contour
    \draw[blue!60, dashed, thick] (1.5, 2.2) -- (3.5, 0.8);

    \node[mylabel] at (3, -3) {等高线越密，梯度越大};
    \node[mylabel] at (3, -3.5) {梯度方向$\perp$等高线};
\end{tikzpicture}
\caption{梯度的几何意义。梯度方向与函数等高线垂直，指向函数值增长最快的方向。负梯度方向指向函数值下降最快的方向。}
\label{fig:gradient}
\end{figure}

\subsection{Jacobian 矩阵}

当输出是多维向量时（如神经网络的某一层输出），推广梯度概念得到 Jacobian 矩阵。设 $\mathbf{f}:\R^n\to\R^m$，其中 $\mathbf{f}(\mathbf{x})=[f_1(\mathbf{x}),\ldots,f_m(\mathbf{x})]^{\top}$，则 Jacobian 矩阵为：

\[
\mathbf{J} = \frac{\partial \mathbf{f}}{\partial \mathbf{x}} =
\begin{bmatrix}
\frac{\partial f_1}{\partial x_1} & \cdots & \frac{\partial f_1}{\partial x_n} \\
\vdots & \ddots & \vdots \\
\frac{\partial f_m}{\partial x_1} & \cdots & \frac{\partial f_m}{\partial x_n}
\end{bmatrix} \in \R^{m\times n}
\]

Jacobian 矩阵的每一行是对应输出分量关于所有输入变量的梯度，每一列是某一输入对所有输出的影响。在反向传播中，Jacobian 矩阵描述了每一层的局部导数。

\begin{mycode}{Python中的梯度计算（数值微分）}
\begin{lstlisting}[language=Python]
import numpy as np

def numerical_gradient(f, x, h=1e-5):
    """用中心差分法计算 f 在 x 处的梯度"""
    grad = np.zeros_like(x)
    it = np.nditer(x, flags=['multi_index'])
    while not it.finished:
        idx = it.multi_index
        x_plus = x.copy()
        x_minus = x.copy()
        x_plus[idx] += h
        x_minus[idx] -= h
        grad[idx] = (f(x_plus) - f(x_minus)) / (2 * h)
        it.iternext()
    return grad

# 测试函数: f(x, y) = x^2 + y^2
def f(x):
    return x[0]**2 + x[1]**2

x0 = np.array([3.0, 4.0])
grad = numerical_gradient(f, x0)
print(f"f(x) = x0^2 + x1^2")
print(f"在 x={x0} 处的数值梯度: {grad}")
print(f"解析梯度 (2x): {2 * x0}")

# 测试: f(x, y) = x^2 * y + sin(y)
def g(x):
    return x[0]**2 * x[1] + np.sin(x[1])

x1 = np.array([2.0, np.pi/4])
grad_g = numerical_gradient(g, x1)
print(f"\n在 x={x1} 处的梯度: {grad_g}")
# 解析: [2xy, x^2 + cos(y)]
analytic = np.array([2*2*np.pi/4, 4 + np.cos(np.pi/4)])
print(f"解析梯度: {analytic}")
\end{lstlisting}
\end{mycode}

\section{链式法则}

\subsection{一元函数的链式法则}

若 $y=f(u)$ 且 $u=g(x)$，则复合函数 $y=f(g(x))$ 的导数为：

\[
\frac{dy}{dx} = \frac{dy}{du} \cdot \frac{du}{dx} = f'(g(x)) \cdot g'(x)
\]

\subsection{多元函数的链式法则}

对于多元复合函数 $z=f(u,v)$ 其中 $u=g(x,y)$、$v=h(x,y)$：

\[
\frac{\partial z}{\partial x} = \frac{\partial z}{\partial u}\frac{\partial u}{\partial x} + \frac{\partial z}{\partial v}\frac{\partial v}{\partial x}
\]

类似地求 $\frac{\partial z}{\partial y}$。这可以推广到任意多个中间变量。

\subsection{反向传播与链式法则}

反向传播（Backpropagation）本质上是链式法则在计算图上的高效实现。考虑一个简单的三层网络：

\[
\mathbf{h}^{(1)} = \sigma(\mathbf{W}^{(1)}\mathbf{x}), \quad
\mathbf{h}^{(2)} = \sigma(\mathbf{W}^{(2)}\mathbf{h}^{(1)}), \quad
\hat{\mathbf{y}} = \mathbf{W}^{(3)}\mathbf{h}^{(2)}
\]

损失函数 $L = \ell(\hat{\mathbf{y}}, \mathbf{y})$。梯度计算如下：

\[
\frac{\partial L}{\partial \mathbf{W}^{(3)}} = \frac{\partial L}{\partial \hat{\mathbf{y}}} \cdot \frac{\partial \hat{\mathbf{y}}}{\partial \mathbf{W}^{(3)}}
\]

\[
\frac{\partial L}{\partial \mathbf{W}^{(2)}} = \frac{\partial L}{\partial \hat{\mathbf{y}}} \cdot
\frac{\partial \hat{\mathbf{y}}}{\partial \mathbf{h}^{(2)}} \cdot
\frac{\partial \mathbf{h}^{(2)}}{\partial \mathbf{W}^{(2)}}
\]

\[
\frac{\partial L}{\partial \mathbf{W}^{(1)}} = \frac{\partial L}{\partial \hat{\mathbf{y}}} \cdot
\frac{\partial \hat{\mathbf{y}}}{\partial \mathbf{h}^{(2)}} \cdot
\frac{\partial \mathbf{h}^{(2)}}{\partial \mathbf{h}^{(1)}} \cdot
\frac{\partial \mathbf{h}^{(1)}}{\partial \mathbf{W}^{(1)}}
\]

注意每一步需要乘以上一层的局部梯度（Jacobian矩阵），这正是"反向"的含义——梯度从输出层逐层传播回输入层。

\begin{figure}[H]
\centering
\begin{tikzpicture}[
    mynode/.style={circle, draw=blue!60, fill=blue!10, minimum size=0.8cm, font=\small},
    myedge/.style={->, >=stealth, thick},
    mygrad/.style={->, >=stealth, thick, red!70, bend right=30},
    mylabel/.style={font=\footnotesize}
]
    % Forward pass nodes
    \node[mynode] (x) at (0, 0) {$\mathbf{x}$};
    \node[mynode] (h1) at (2.5, 0) {$\mathbf{h}^{(1)}$};
    \node[mynode] (h2) at (5, 0) {$\mathbf{h}^{(2)}$};
    \node[mynode] (y) at (7.5, 0) {$\hat{\mathbf{y}}$};
    \node[mynode, fill=red!10, draw=red!60] (L) at (10, 0) {$L$};

    % Forward edges
    \draw[myedge, blue!70] (x) -- (h1) node[midway, above] {\tiny $\mathbf{W}^{(1)}$};
    \draw[myedge, blue!70] (h1) -- (h2) node[midway, above] {\tiny $\mathbf{W}^{(2)}$};
    \draw[myedge, blue!70] (h2) -- (y) node[midway, above] {\tiny $\mathbf{W}^{(3)}$};
    \draw[myedge, blue!70] (y) -- (L);

    % Gradient edges (backward)
    \draw[mygrad] (L) to[bend right=40] node[below] {\tiny $\frac{\partial L}{\partial\mathbf{W}^{(3)}}$} (y);
    \draw[mygrad] (L) to[bend right=50] node[below=0.7cm] {\tiny $\frac{\partial L}{\partial\mathbf{W}^{(2)}}$} (h2);
    \draw[mygrad] (L) to[bend right=60] node[below=1.4cm] {\tiny $\frac{\partial L}{\partial\mathbf{W}^{(1)}}$} (h1);

    \node[mylabel] at (5, -2.2) {红色：反向传播路径};
    \node[mylabel] at (5, -2.7) {蓝色：前向传播路径};
\end{tikzpicture}
\caption{反向传播示意图。前向传播（蓝色）计算每一层的输出，反向传播（红色）从损失函数出发，利用链式法则逐层计算梯度。}
\label{fig:backprop}
\end{figure}

\begin{mycode}{Python中的简单反向传播实现}
\begin{lstlisting}[language=Python]
import numpy as np

np.random.seed(42)

# 简单三层网络
n_input = 4; n_hidden1 = 3; n_hidden2 = 2; n_output = 1

W1 = np.random.randn(n_hidden1, n_input) * 0.1
W2 = np.random.randn(n_hidden2, n_hidden1) * 0.1
W3 = np.random.randn(n_output, n_hidden2) * 0.1

x = np.random.randn(n_input, 1)
y_true = np.array([[1.0]])

def sigmoid(z):
    return 1 / (1 + np.exp(-z))

# 前向传播
z1 = W1 @ x;            a1 = sigmoid(z1)
z2 = W2 @ a1;           a2 = sigmoid(z2)
z3 = W3 @ a2;           a3 = z3  # 线性输出

# 损失 (MSE)
loss = 0.5 * np.sum((a3 - y_true) ** 2)
print(f"前向损失: {loss:.6f}")

# 反向传播
delta3 = a3 - y_true                         # dL/da3
dW3 = delta3 @ a2.T                          # dL/dW3

delta2 = (W3.T @ delta3) * (a2 * (1 - a2))  # dL/da2 -> dL/dz2
dW2 = delta2 @ a1.T                          # dL/dW2

delta1 = (W2.T @ delta2) * (a1 * (1 - a1))  # dL/da1 -> dL/dz1
dW1 = delta1 @ x.T                           # dL/dW1

print(f"dL/dW3 形状: {dW3.shape}")
print(f"dL/dW2 形状: {dW2.shape}")
print(f"dL/dW1 形状: {dW1.shape}")

# 手动梯度下降一步
lr = 0.01
W3 -= lr * dW3
W2 -= lr * dW2
W1 -= lr * dW1
print(f"一步更新后的权重大小: {np.sum(np.abs(W1)):.4f}")
\end{lstlisting}
\end{mycode}

\section{Hessian 矩阵}

\begin{mydefinition}{Hessian 矩阵}
对于二阶可导的函数 $f:\R^n\to\R$，其 Hessian 矩阵 $\mathbf{H}$ 是由所有二阶偏导数组成的 $n\times n$ 矩阵：
\[
\mathbf{H}_{ij} = \frac{\partial^2 f}{\partial x_i \partial x_j}
\]
\end{mydefinition}

Hessian 矩阵是对称矩阵（当二阶混合偏导数连续时，即 $\frac{\partial^2 f}{\partial x_i\partial x_j}=\frac{\partial^2 f}{\partial x_j\partial x_i}$）。

Hessian 矩阵在深度学习中的作用：
\begin{itemize}
    \item \textbf{驻点判断}：若 $\nabla f(\mathbf{x}^*)=0$ 且 $\mathbf{H}(\mathbf{x}^*)$ 正定，则 $\mathbf{x}^*$ 是局部极小点。
    \item \textbf{牛顿法}：利用 Hessian 信息实现二阶优化：$\mathbf{x}_{t+1}=\mathbf{x}_t-\mathbf{H}^{-1}\nabla f(\mathbf{x}_t)$。
    \item \textbf{曲率信息}：Hessian 的特征值反映损失函数在不同方向上的弯曲程度。最大特征值对应最陡峭的方向，最小特征值对应最平坦的方向。
    \item \textbf{鞍点分析}：在高维优化中，Hessian 同时有正负特征值意味着存在鞍点，这是深度学习优化的主要挑战之一。
\end{itemize}

由于深度学习参数数量巨大，直接计算和存储完整的 Hessian 矩阵不现实，因此实际中常用近似方法（如 BFGS、L-BFGS、Hessian-Free优化）。

\section{Taylor 展开}

Taylor 展开是将函数在某点附近用多项式逼近的方法。它是理解优化算法收敛性的理论基础。

\subsection{一元函数的 Taylor 展开}

若 $f$ 在 $x_0$ 处 $n+1$ 阶可导，则在 $x_0$ 附近：

\[
f(x) = \sum_{k=0}^{n} \frac{f^{(k)}(x_0)}{k!}(x-x_0)^k + R_n(x)
\]

其中 $R_n(x)=\frac{f^{(n+1)}(\xi)}{(n+1)!}(x-x_0)^{n+1}$ 为 Lagrange 余项，$\xi$ 在 $x$ 和 $x_0$ 之间。

常用的一阶和二阶 Taylor 展开：
\begin{itemize}
    \item 一阶：$f(x)\approx f(x_0)+f'(x_0)(x-x_0)$（线性近似）
    \item 二阶：$f(x)\approx f(x_0)+f'(x_0)(x-x_0)+\frac{1}{2}f''(x_0)(x-x_0)^2$（二次近似）
\end{itemize}

\subsection{多元函数的 Taylor 展开}

对于多元函数 $f(\mathbf{x})$，在 $\mathbf{x}_0$ 附近的二阶 Taylor 展开为：

\[
f(\mathbf{x}) \approx f(\mathbf{x}_0) + \nabla f(\mathbf{x}_0)^{\top}(\mathbf{x}-\mathbf{x}_0) + \frac{1}{2}(\mathbf{x}-\mathbf{x}_0)^{\top}\mathbf{H}(\mathbf{x}_0)(\mathbf{x}-\mathbf{x}_0)
\]

这一展开式在分析梯度下降的收敛性时非常重要。例如，用二阶 Taylor 展开分析学习率对收敛性的影响：设 $\mathbf{x}_{t+1}=\mathbf{x}_t-\eta\nabla f(\mathbf{x}_t)$，代入上式得：

\[
f(\mathbf{x}_{t+1}) \approx f(\mathbf{x}_t) - \eta\|\nabla f(\mathbf{x}_t)\|^2 + \frac{\eta^2}{2}\nabla f(\mathbf{x}_t)^{\top}\mathbf{H}\nabla f(\mathbf{x}_t)
\]

为使损失下降，需要 $\eta < \frac{2\|\nabla f\|^2}{\nabla f^{\top}\mathbf{H}\nabla f}$。

\begin{mycode}{Python中的Taylor展开验证}
\begin{lstlisting}[language=Python]
import numpy as np
import matplotlib.pyplot as plt

def f(x):
    return np.sin(x) + 0.1 * x**2

def f_prime(x):
    return np.cos(x) + 0.2 * x

def f_double_prime(x):
    return -np.sin(x) + 0.2

x0 = 1.0  # 展开点

# 计算Taylor展开
x = np.linspace(-2, 4, 200)
fx = f(x)

# 一阶Taylor
taylor1 = f(x0) + f_prime(x0) * (x - x0)

# 二阶Taylor
taylor2 = taylor1 + 0.5 * f_double_prime(x0) * (x - x0)**2

# 三阶Taylor
f_triple = lambda x: -np.cos(x)
taylor3 = taylor2 + (f_triple(x0) / 6) * (x - x0)**3

print(f"展开点 x0 = {x0}")
print(f"f(x0) = {f(x0):.4f}")
print(f"在 x=0.5 处:")
print(f"  真实值: {f(0.5):.4f}")
print(f"  一阶近似: {f(x0) + f_prime(x0)*(0.5-x0):.4f}")
print(f"  二阶近似: {f(x0) + f_prime(x0)*(0.5-x0) + 0.5*f_double_prime(x0)*(0.5-x0)**2:.4f}")
\end{lstlisting}
\end{mycode}

\begin{figure}[H]
\centering
\begin{tikzpicture}[
    mycurve/.style={thick, smooth},
    mypoint/.style={circle, fill=red, inner sep=2pt}
]
    % Axes
    \draw[->, >=stealth] (-3, 0) -- (5, 0) node[right] {$x$};
    \draw[->, >=stealth] (0, -2) -- (0, 3) node[above] {$f(x)$};

    % True function (like sin(x))
    \draw[mycurve, blue!70, domain=-2.5:4.5, samples=100]
        plot (\x, {sin(\x r) + 0.1*\x*\x});

    % First order Taylor at x0=1
    \draw[mycurve, red!70, domain=-1:3]
        plot (\x, {sin(1 r) + 0.1 + (cos(1 r) + 0.2)*(\x - 1)});

    % Second order Taylor
    \draw[mycurve, green!50!black, domain=-1:3, samples=100]
        plot (\x, {sin(1 r)+0.1+(cos(1 r)+0.2)*(\x-1)+0.5*(-sin(1 r)+0.2)*(\x-1)*(\x-1)});

    % Expansion point
    \node[mypoint] at (1, {sin(1 r)+0.1}) {};
    \node[font=\small] at (1.2, {sin(1 r)+0.1+0.3}) {$x_0$};

    % Legend
    \node[font=\small, blue!70, right] at (4, 2) {原函数};
    \node[font=\small, red!70, right] at (4, 1.5) {一阶Taylor};
    \node[font=\small, green!50!black, right] at (4, 1) {二阶Taylor};
\end{tikzpicture}
\caption{函数在一阶和二阶 Taylor 展开下的近似。在展开点附近，二阶近似比一阶近似更精确，但远离展开点时近似误差会增大。}
\label{fig:taylor}
\end{figure}

\section{凸优化基础}

\subsection{凸集与凸函数}

\begin{mydefinition}{凸集}
集合 $C\subseteq\R^n$ 是凸集，若对任意 $\mathbf{x},\mathbf{y}\in C$ 和任意 $\lambda\in[0,1]$，有 $\lambda\mathbf{x}+(1-\lambda)\mathbf{y}\in C$。即连接集合中任意两点的线段完全包含在集合中。
\end{mydefinition}

\begin{mydefinition}{凸函数}
定义在凸集 $C$ 上的函数 $f$ 是凸函数，若对任意 $\mathbf{x},\mathbf{y}\in C$ 和任意 $\lambda\in[0,1]$，有：
\[
f(\lambda\mathbf{x}+(1-\lambda)\mathbf{y}) \leq \lambda f(\mathbf{x}) + (1-\lambda)f(\mathbf{y})
\]
即函数图像在连接任意两点的线段下方。
\end{mydefinition}

若不等式严格成立（$\mathbf{x}\neq\mathbf{y}$ 且 $\lambda\in(0,1)$ 时），则称 $f$ 为严格凸函数。

\subsection{凸函数的判定条件}

\begin{myTheorem}{凸函数的判定}
设 $f:\R^n\to\R$ 二阶可导，则：
\begin{itemize}
    \item $f$ 是凸函数当且仅当其 Hessian 矩阵处处半正定：$\mathbf{H}(\mathbf{x})\succeq 0$，$\forall\mathbf{x}$
    \item 若 $\mathbf{H}(\mathbf{x})\succ 0$（正定），则 $f$ 是严格凸函数
\end{itemize}
\end{myTheorem}

对于一元函数，条件简化为 $f''(x)\geq 0$（处处成立）。

常见的凸函数：
\begin{itemize}
    \item $f(x)=x^2$：严格凸
    \item $f(x)=e^x$：严格凸
    \item $f(x)=|x|$：凸但不严格
    \item $f(x)=x\log x$（$x>0$）：严格凸
    \item MSE 损失 $f(\mathbf{w})=\|\mathbf{X}\mathbf{w}-\mathbf{y}\|_2^2$：凸
    \item 带有 $L^2$ 正则化的线性回归损失：严格凸
\end{itemize}

\subsection{凸优化的性质}

凸优化问题具有优良的性质：

\begin{enumerate}
    \item \textbf{局部极小即全局极小}：对凸函数，任何局部极小点都是全局极小点。
    \item \textbf{驻点即极小点}：若 $\nabla f(\mathbf{x}^*)=0$，则 $\mathbf{x}^*$ 是全局极小点。
    \item \textbf{梯度下降收敛到全局最优}：对于凸函数，梯度下降可以保证收敛到全局最优解。
\end{enumerate}

然而，深度学习的损失函数通常是高度非凸的，存在大量局部极小点和鞍点。但实践表明，SGD 等方法在深度网络中表现良好，部分原因在于：
\begin{itemize}
    \item 高维空间中的鞍点比局部极小点更常见
    \item SGD的随机噪声有助于逃离鞍点
    \item 过参数化的网络往往使优化问题变得"更容易"
\end{itemize}

\begin{mycode}{Python中的凸函数与非凸函数对比}
\begin{lstlisting}[language=Python]
import numpy as np

# 凸函数: f(x) = x^2
def convex_f(x):
    return x**2

# 非凸函数: f(x) = sin(x) + 0.1 * x^2
def nonconvex_f(x):
    return np.sin(3*x) + 0.1 * x**2

# 验证凸性: 检查 Jensen 不等式
def check_convexity(f, x1, x2, n_points=100):
    violations = 0
    for lam in np.linspace(0, 1, n_points):
        mid = lam * x1 + (1 - lam) * x2
        f_mid = f(mid)
        bound = lam * f(x1) + (1 - lam) * f(x2)
        if f_mid > bound + 1e-10:
            violations += 1
    return violations

v1 = check_convexity(convex_f, -2, 3)
v2 = check_convexity(nonconvex_f, -2, 3)
print(f"f(x)=x^2 违反凸性次数: {v1} (应为0)")
print(f"f(x)=sin(3x)+0.1x^2 违反凸性次数: {v2} (预期>0)")
\end{lstlisting}
\end{mycode}

\section{约束优化与 Lagrange 乘子法}

在深度学习中，有时需要在约束条件下优化目标函数。Lagrange 乘子法是处理等式约束优化的基本方法。

\begin{myTheorem}{Lagrange 乘子法}
考虑约束优化问题：$\min_{\mathbf{x}} f(\mathbf{x})$，满足 $g(\mathbf{x})=0$。构造 Lagrange 函数：
\[
\mathcal{L}(\mathbf{x},\lambda) = f(\mathbf{x}) + \lambda g(\mathbf{x})
\]
则最优解满足 $\nabla_{\mathbf{x}}\mathcal{L}=0$ 和 $\nabla_{\lambda}\mathcal{L}=g(\mathbf{x})=0$。
\end{myTheorem}

对于不等式约束 $g(\mathbf{x})\leq 0$，使用 KKT 条件（Karush-Kuhn-Tucker conditions）进行推广。SVM（支持向量机）的对偶推导就是 Lagrange 乘子法的典型应用。

\section{本章小结}

本章介绍了深度学习所需的微积分基础知识：
\begin{itemize}
    \item 导数和偏导数描述了函数的变化率，梯度指向函数增长最快的方向
    \item 链式法则是反向传播算法的数学基础
    \item Hessian 矩阵提供了二阶曲率信息，用于判断驻点类型
    \item Taylor 展开将复杂函数局部近似为多项式
    \item 凸函数保证局部极小即全局极小
    \item 在深度学习中，梯度下降、牛顿法和各种优化器都建立在微积分理论之上
\end{itemize}

这些概念将在后续章节中反复出现。下一章将介绍概率论与统计学基础，为理解损失函数的概率解释和贝叶斯方法做准备。
